\section{Autoencoders geométricos: desafíos técnicos clave}

\subsection{Problema: los VAE de imagen no son adecuados para datos geométricos}

B3D utiliza modelos de difusión latentes para generar simultáneamente imágenes RGB y nubes de puntos XYZ, lo que requiere autoencodificar ambas modalidades. Sin embargo, Henzler descubrió que \textbf{el uso directo de VAE de imagen preentrenados para codificar y decodificar nubes de puntos genera resultados extremadamente pobres}, y el efecto empeora a medida que aumenta la escala de la escena.

Análisis de las causas:
\begin{itemize}
    \item \textbf{Incompatibilidad de funciones de pérdida}: la pérdida perceptual, común en la compresión de imágenes, carece de sentido para datos geométricos.
    \item \textbf{Incompatibilidad del rango de los datos}: los valores de los píxeles de una imagen siempre se encuentran en el intervalo $[0, 1]$, mientras que el rango de valores de las coordenadas geométricas es completamente diferente y varía según la escala de la escena.
\end{itemize}

\begin{warningbox}{Trampa de usar VAE de imagen para codificación geométrica}
Aplicar directamente un VAE diseñado para imágenes (entrenado con pérdidas perceptuales, GAN, etc.) a datos geométricos (nubes de puntos XYZ) no es viable. Los datos geométricos difieren esencialmente de los datos de imagen en términos de distribución estadística, rango de valores y significado semántico. Es imperativo diseñar un autoencoder geométrico especializado.
\end{warningbox}

\subsection{Arquitectura del decodificador: Transformer vs. Convoluciones}

Otro hallazgo interesante se refiere a la elección de la arquitectura del decodificador:

\begin{itemize}
    \item Desde una perspectiva cuantitativa, la calidad de reconstrucción geométrica de los decodificadores basados en convoluciones y en Transformer es comparable.
    \item Sin embargo, la \textbf{observación cualitativa} revela diferencias clave: los decodificadores basados en convoluciones generan artefactos visualmente molestos, por ejemplo, incapaces de mantener la linealidad de las líneas rectas.
    \item Por lo tanto, B3D eligió un \textbf{decodificador Transformer} para procesar los datos XYZ.
\end{itemize}

\begin{warningbox}{Diferencia entre indicadores cuantitativos y evaluación cualitativa}
Este es un caso típico en el que los indicadores cuantitativos (como el MSE) pueden no reflejar completamente la calidad visual real. Aunque ambos decodificadores muestran un rendimiento similar en métricas numéricas, la salida del decodificador basado en convoluciones presenta artefactos estructurales (como líneas rectas curvas), lo cual es inaceptable en aplicaciones prácticas. Esto nos recuerda que, al evaluar la calidad de reconstrucción geométrica, no debemos depender únicamente de las métricas tradicionales.
\end{warningbox}

\subsection{Resumen de la sección}

El autoencoder geométrico es un componente clave de los modelos de difusión latente 3D, pero su diseño enfrenta desafíos distintos a los de los autoencoders de imagen. Las diferencias en la distribución de los datos, la selección de la función de pérdida y la arquitectura del decodificador requieren una reconsideración adaptada a las características específicas de los datos geométricos.
